% Incorporación aditiva para el Artículo I; no modifica la fuente anterior.
% RESULTADO_RECUPERADO: integral, capítulo 22, propietario extenso
% colaboracion/partes_i_ii/source/base_residencias/
% capitulo_22_c20_lineas_2613_3270.tex, líneas 176--576.
% Equivalencia 2/3: public/residencias/capitulo_22_voa_leech_moonshine.tex,
% líneas 55--127. Orden dos y Fricke: hmt/15d_voa_moonshine_orden_dos_rev11.tex.
% Prerrequisitos residentes en I: código C_W, matriz A_W, N(A_2^12),
% origen o_*, vector v_alpha de norma54, N_alpha reconocido como Leech,
% VOA reticular y presentación FLM de las secciones precedentes.
% Bibliografía nueva: chenlamshimakura2016, abelamyamada2017,
% kmconwaynorton1979 (bibitems listos para incorporar al final de este archivo).

\subsection{Orientación ternaria del vecino de Leech}
\label{km:orden-tres-reticular}

La presentación de orden dos construida anteriormente utiliza la negación
del retículo de Leech. El mismo retículo, obtenido de la incidencia marcada
por \(K\) y por el soporte dodecafásico de \(\alpha\), admite una segunda
presentación compatible con el carácter ternario de la construcción.
Para hacer explícita esta continuidad se conserva el pegado, se construye
su isometría de orden tres y se comprueba que ésta preserva el vecino ya
seleccionado. El paso no requiere alterar \(K\), su lectura racional ni
la definición previa de \(\alpha\). La proposición
\ref{km:estabilidad-vecino} verifica la isometría reticular; la proposición
\ref{km:traza-tipo} y el teorema~\ref{km:moonshine-tres} precisan la construcción
de orden tres y las hipótesis de los reconocimientos clásicos utilizados.

Trabajamos en la base de raíces simples de cada factor \(A_2\), con
\begin{equation}
 B_2=\begin{pmatrix}2&-1\\-1&2\end{pmatrix},\qquad
 \mathsf C=\begin{pmatrix}0&-1\\1&-1\end{pmatrix},\qquad
 \mathsf R=\begin{pmatrix}0&1\\1&0\end{pmatrix},\qquad
 \rho=\begin{pmatrix}1\\1\end{pmatrix}.
 \label{km:coxeter-reflexion}
\end{equation}
La multiplicación da
\[
 \mathsf C^3=I_2,\quad
 \mathsf C^2+\mathsf C+I_2=0,\quad
 \mathsf C^{\mathsf T}B_2\mathsf C=B_2,\quad
 \mathsf R\mathsf C\mathsf R=\mathsf C^{-1},\quad
 \mathsf R^{\mathsf T}B_2\mathsf R=B_2.
\]
Estas identidades fijan tanto la métrica como la orientación de la acción.
En el discriminante \(A_2^*/A_2\cong\mathbb F_3\) tomamos como
generador la clase de \(\varpi=(2,1)^{\mathsf T}/3\). Se cumple
\(\mathsf C\varpi-\varpi=(-1,0)^{\mathsf T}\), por lo que
\(\mathsf C\) actúa trivialmente sobre el discriminante; la reflexión
\(\mathsf R\) cambia su signo.

La matriz de Paley--Witt ya construida produce un elemento de soporte
completo mediante la siguiente operación:
\begin{equation}
 \begin{gathered}
 q_*=(1,1,1,1,1,1),\qquad q_*A_W=(2,1,1,1,1,1),\\
 c_*=(q_*,q_*A_W)
 =(1,1,1,1,1,1,2,1,1,1,1,1)\in C_W.
 \end{gathered}
 \label{km:palabra-total}
\end{equation}
Cada coordenada de \(c_*\) es no nula. Para las doce posiciones
definimos los marcos y la isometría
\begin{equation}
 P_b(c_*)=
 \begin{cases}\mathsf R,&(c_*)_b=1,\\I_2,&(c_*)_b=2,\end{cases}
 \qquad
 g_c=\bigoplus_{b=1}^{12}P_b(c_*)\mathsf C P_b(c_*)^{-1}.
 \label{km:isometria-ternaria}
\end{equation}
El símbolo \(P_b(c_*)\) designa aquí un marco invertible de un factor
\(A_2\), no el proyector sobre el isótipo tridimensional \(P_3\)
de la sección~\ref{exc:k-direccion}.

\begin{proposition}[Isometría de orden tres sobre el vecino marcado]
\label{km:estabilidad-vecino}
El operador \(g_c\) es una isometría de orden tres sin vectores fijos
no nulos. Preserva el retículo de Niemeier \(N=N(C_W)\) y el vecino
\(N_\alpha\) construido en \eqref{exc:vecino}.
\end{proposition}
\begin{proof}
Cada bloque de \(g_c\) es \(\mathsf C\) o \(\mathsf C^{-1}\);
su polinomio mínimo es \(X^2+X+1\). Esto prueba el orden y la
ausencia de vectores fijos. Ambos bloques actúan trivialmente sobre el
discriminante, de modo que conservan cada clase de pegado y, por tanto,
el retículo \(N\).

Escribamos \(v=v_\alpha=4\rho_{o_*}+\sum_{b\ne o_*}\rho_b\),
con los marcos y el origen incidencial ya fijados. Sus coeficientes
radiales son \(a_{o_*}=4\) y \(a_b=1\) para \(b\ne o_*\).
Todos satisfacen \(a_b\equiv1\pmod3\). En un factor,
\[
 (\mathsf C-I_2)\rho=(-2,-1)^{\mathsf T},\qquad
 (\mathsf C^{-1}-I_2)\rho=(-1,-2)^{\mathsf T}.
\]
Después de dividir por tres, sus clases discriminantes son,
respectivamente, \(2\) y \(1\). En consecuencia, los vectores
\begin{equation}
 y=\frac{g_cv-v}{3},\qquad z=\frac{g_c^{-1}v-v}{3}
 \label{km:desplazamientos}
\end{equation}
tienen palabras discriminantes \(c_*\) y \(-c_*\), ambas en
\(C_W\). Por ello \(y,z\in N\). El producto local es
\[
 \left\langle
 \frac{(\mathsf C^{\pm1}-I_2)a_b\rho}{3},a_b\rho
 \right\rangle=-a_b^2,
\]
y su suma da
\(\langle y,v\rangle=\langle z,v\rangle=-(16+11)=-27\).
Así, \(y,z\in N_0=\{x\in N:\langle x,v\rangle\equiv0\pmod3\}\).
Para \(x\in N_0\),
\[
 \langle g_cx,v\rangle
 =\langle x,g_c^{-1}v\rangle
 =\langle x,v\rangle+3\langle x,z\rangle\equiv0\pmod3.
\]
El argumento aplicado a \(g_c^{-1}\) prueba \(g_cN_0=N_0\).
Finalmente, \(g_c(v/3)=v/3+y\) conserva la clase que genera
\(N_\alpha=N_0+\mathbb Z(v/3)\). Se obtiene
\(g_cN_\alpha=N_\alpha\).
\end{proof}

La prueba conserva más que la dimensión del retículo: utiliza el código
de pegado, su orientación, el origen de la bandera y la congruencia del
vector radial. Estos datos explican por qué la segunda presentación
permanece vinculada al registro dodecafásico del que partió la construcción.

\subsection{Traza ternaria, peso torcido y extensión de orden tres}
\label{km:orbifold-tres}

Sea \(\Lambda=N_\alpha\) y sea \(\widehat g_c\) el levantamiento
estándar de \(g_c\) al álgebra reticular
\(V_\Lambda=M(1)\otimes\mathbb C_\varepsilon[\Lambda]\), con la
extensión central y el cociclo ya declarados. El orden impar permite
elegir este levantamiento de orden tres. Cada factor complejo \(A_2\)
aporta los autovalores \(\zeta_3\) y \(\zeta_3^2\), donde
\(\zeta_3^2+\zeta_3+1=0\).

\begin{proposition}[Traza graduada y tipo del levantamiento]
\label{km:traza-tipo}
Para \(j=1,2\),
\begin{equation}
 \operatorname{Tr}_{V_\Lambda}
       \bigl(\widehat g_c^{\,j}q^{L_0-1}\bigr)
 =t_3(\tau)
 :=q^{-1}\prod_{n\ge1}(1+q^n+q^{2n})^{-12}.
 \label{km:tres-traza}
\end{equation}
El peso mínimo de los módulos torcidos no triviales es \(4/3\), y
el levantamiento es de tipo cero.
\end{proposition}
\begin{proof}
En la parte reticular de la traza sólo contribuye el vector cero,
porque \(g_c\) y \(g_c^2\) carecen de vectores fijos no nulos.
En cada grado oscilador, la traza de potencias simétricas es el inverso
del determinante \(I-g_c^jq^n\). Cada factor \(A_2\) contribuye
\(1+q^n+q^{2n}\), lo que demuestra \eqref{km:tres-traza}.

La fórmula del peso del vacío torcido para un automorfismo reticular de
orden tres, aplicada a los dos espacios propios de dimensión doce, da
\[
 \rho_{\widehat g_c}
 =\frac{1}{4\cdot3^2}
   \bigl(1\cdot2\cdot12+2\cdot1\cdot12\bigr)=\frac43.
\]
El tipo es la clase de \(3^2\rho_{\widehat g_c}=12\) módulo
tres, que es cero. La fórmula de peso es un resultado de la teoría de
módulos torcidos; las multiplicidades y la isometría que la alimentan
se han construido aquí.
\end{proof}

\begin{theorem}[Segunda presentación del álgebra Moonshine]
\label{km:moonshine-tres}
La extensión cíclica de tipo cero
\begin{equation}
 V_{(3)}=\operatorname{Orb}_{\mathbb Z/3}
                 (V_\Lambda,\widehat g_c)
 \label{km:extension-orden3}
\end{equation}
es isomorfa al álgebra de operadores de vértice \(V^\natural\).
\end{theorem}
\begin{proof}
El teorema \ref{exc:teorema-leech} proporciona un retículo positivo,
par, unimodular, de rango veinticuatro y sin raíces. La
Proposición~\ref{km:estabilidad-vecino} construye sobre ese mismo
retículo una isometría de orden tres sin vectores fijos; la
Proposición~\ref{km:traza-tipo} verifica el levantamiento y su tipo.
Se aplica entonces el teorema de Chen--Lam--Shimakura sobre el
orbifold de orden tres del álgebra reticular de Leech
\cite[Teorema~1.1]{chenlamshimakura2016}. La identificación también queda comprendida
en el tratamiento uniforme de Abe--Lam--Yamada
\cite[Teorema~4.4 y secciones~2--4]{abelamyamada2017}. Esta aplicación es posterior a la construcción
de los datos reticulares; la igualdad de caracteres no sustituye las
hipótesis del teorema utilizado.
\end{proof}

Para precisar la simetría dual, escribamos
\(V_{(3)}=W_0\oplus W_1\oplus W_2\), donde \(W_i\) es el sector
integral seleccionado en el módulo \(\widehat g_c^{\,i}\)-torcido.
El producto satisface
\(Y(W_i,z)W_j\subset W_{i+j\bmod3}((z))\). Por ello
\begin{equation}
 g^\vee|_{W_i}=\zeta_3^i I
 \label{km:dual-ciclico}
\end{equation}
define un automorfismo de orden tres. En efecto, el factor sobre un
producto es \(\zeta_3^{i+j}=\zeta_3^i\zeta_3^j\), y el vacío y el
vector conforme pertenecen a \(W_0\). La identificación clásica de
este automorfismo en el Monstruo es la clase \(3B\), con serie
\(T_{3B}=t_3+12\) \cite{kmconwaynorton1979,borcherds}.
Esta normalización puede comprobarse en la construcción. El carácter
reticular es \(\operatorname{ch}V_\Lambda=J+24\): su polo inicial
es \(q^{-1}\), su espacio de peso uno tiene dimensión 24 y el
carácter modular de peso cero queda determinado por esos datos.
Si \(F\) es el carácter del subespacio fijo y \(H\) el de cada
sector torcido integral, la proyección sobre invariantes y la conjugación
de los dos sectores dan
\[
 F=\frac{J+24+2t_3}{3},\qquad F+2H=J,
 \qquad \operatorname{Tr}g^\vee q^{L_0-1}=F-H=t_3+12.
\]
Su dominio es el álgebra extendida, mientras
\(\widehat g_c\) actúa sobre el álgebra reticular anterior.

\subsection{Equivalencia de las presentaciones de órdenes dos y tres}
\label{km:equivalencia-dos-tres}

Denotemos por
\[
 V_{(2)}=(V_\Lambda)^+\oplus(V_\Lambda^T)^+
\]
la presentación FLM de \eqref{exc:flm}. Ambas presentaciones conservan
el mismo vecino de Leech como antecedente. Sus extensiones emplean
automorfismos diferentes y productos de vértice determinados mediante
los respectivos datos de levantamiento.

\begin{theorem}[Equivalencia graduada y transporte de automorfismos]
\label{km:equivalencia-voa}
Fijados isomorfismos de álgebras de operadores de vértice
\[
 \phi_2:V_{(2)}\xrightarrow{\ \cong\ }V^\natural,
 \qquad
 \phi_3:V_{(3)}\xrightarrow{\ \cong\ }V^\natural,
\]
la composición
\begin{equation}
 \Phi_{23}=\phi_2^{-1}\phi_3:
 V_{(3)}\xrightarrow{\ \cong\ }V_{(2)}
 \label{km:Phi23}
\end{equation}
conserva el vacío, el vector conforme, la gradación y el producto de
vértice. La conjugación por \(\Phi_{23}\) identifica los grupos de
automorfismos y conserva sus trazas graduadas.
\end{theorem}
\begin{proof}
La composición de los dos isomorfismos preserva las estructuras
declaradas. En particular,
\[
 \Phi_{23}Y_{(3)}(a,z)b
 =Y_{(2)}(\Phi_{23}a,z)\Phi_{23}b,
 \qquad
 \Phi_{23}L_0^{(3)}=L_0^{(2)}\Phi_{23}.
\]
Si \(h\) es un automorfismo de \(V_{(3)}\), entonces
\(h' =\Phi_{23}h\Phi_{23}^{-1}\) es un automorfismo de
\(V_{(2)}\). La igualdad de trazas en cada componente homogénea
finito-dimensional da
\[
 \operatorname{Tr}_{V_{(3)}}h q^{L_0-1}
 =\operatorname{Tr}_{V_{(2)}}h' q^{L_0-1}.
\]
\end{proof}

La elección de \(\phi_2\) y \(\phi_3\) forma parte del enunciado;
\(\Phi_{23}\) no se presenta como una identificación canónica
independiente de esas elecciones. Tampoco transforma un automorfismo
de orden tres en la involución de orden dos: la conjugación conserva
el orden. Su contenido es que dos extensiones distintas realizan la
misma estructura graduada, con transporte explícito de sus productos,
representaciones y trazas.

\subsection{Funciones eta, involución de Fricke y normalización del vacío}
\label{km:fricke}

Para \(\operatorname{Im}\tau>0\), las trazas de los dos automorfismos
reticulares admiten las expresiones
\begin{equation}
 \eta(\tau)=q^{1/24}\prod_{n\ge1}(1-q^n),\qquad
 t_2=\left(\frac{\eta(\tau)}{\eta(2\tau)}\right)^{24},\qquad
 t_3=\left(\frac{\eta(\tau)}{\eta(3\tau)}\right)^{12}.
 \label{km:eta-trazas}
\end{equation}
Se deducen directamente de las factorizaciones
\(1-q^{2n}=(1-q^n)(1+q^n)\) y
\(1-q^{3n}=(1-q^n)(1+q^n+q^{2n})\). El orden en \(q\) del
cociente de nivel tres, antes de elevarlo a la potencia doce, es
\[
 \frac1{24}-\frac3{24}=-\frac1{12};
 \qquad 12\left(-\frac1{12}\right)=-1.
\]
Aquí \(-1/12\) designa la diferencia exacta de los exponentes
iniciales de dos funciones eta. La operación está especificada y no
requiere interpretar una suma divergente como una suma convergente.
La igualdad de este escalar con un defecto normalizado de otra
realización no identifica los operadores que lo producen: cada
realización conserva su dominio y su operación.

\begin{proposition}[Involución de Fricke sobre la traza de orden dos]
\label{km:fricke-prop}
Sea \(w_2\tau=-1/(2\tau)\). Entonces
\begin{equation}
 t_2(w_2\tau)=\frac{2^{12}}{t_2(\tau)},\qquad
 T_{2A}=t_2+24+\frac{2^{12}}{t_2}
 \label{km:fricke-t2}
\end{equation}
y \(T_{2A}(w_2\tau)=T_{2A}(\tau)\).
\end{proposition}
\begin{proof}
La ley de transformación
\(\eta(-1/\tau)=(-i\tau)^{1/2}\eta(\tau)\) da
\[
 \frac{\eta(-1/(2\tau))}{\eta(-1/\tau)}
 =\sqrt2\,\frac{\eta(2\tau)}{\eta(\tau)}.
\]
La potencia veinticuatro produce la primera identidad. La sustitución
intercambia los dos términos variables de \(T_{2A}\) y deja fijo
el término constante. El reconocimiento de esta expresión como serie
de McKay--Thompson de la clase \(2A\) corresponde a la
simetrización de nivel dos de Conway--Norton
\cite[\S\S4 y 6, tabla~3]{kmconwaynorton1979} y al teorema
de Moonshine \cite{borcherds}.
\end{proof}

Las uniformizaciones modulares de niveles dos y tres, en las
normalizaciones anteriores, son
\begin{align}
 j&=\frac{(t_2+256)^3}{t_2^2},
 \label{km:j-nivel2}\\
 J=j-744
 &=t_3+12+\frac{10\cdot3^9}{t_3}
       +\frac{4\cdot3^{14}}{t_3^2}
       +\frac{3^{18}}{t_3^3}.
 \label{km:j-nivel3}
\end{align}
Se utilizan como identidades modulares clásicas sobre las trazas ya
construidas, no como reglas para seleccionar el código o el vecino.
La segunda también se escribe
\(j=(t_3+27)(t_3+243)^3/t_3^3\).
El desarrollo directo del producto de \eqref{km:tres-traza} comienza por
\[
 t_3=q^{-1}-12+54q-76q^2-243q^3+1188q^4+O(q^5).
\]
En particular, \([q]J=54+10\cdot3^9=196884\); ésta es una
consecuencia de las trazas y la uniformización. La construcción de
\(V^\natural\) conserva su prueba reticular y de orbifold,
independientemente de esta comprobación del primer coeficiente.

Las tres operaciones que comparecen en esta prolongación tienen
dominios distintos. La simetría dual \(g^\vee\) actúa sobre sectores
del álgebra de vértices; Fricke actúa sobre el parámetro modular y
su función; la dualidad de radio intercambia momento y enrollamiento
en el sector reticular compacto. Sus fórmulas de reciprocidad pueden
coordinarse mediante los mapas correspondientes, pero el nombre
«dualidad» no sustituye esos mapas. Esta distinción permite continuar
la construcción hacia las pantallas dimensionales manteniendo la
incidencia, la gradación y el origen de cada operación.

% BIBITEMS NUEVOS: copiar al entorno thebibliography del artículo.
% Verificación primaria: arXiv 1606.05961v2 (teorema 1.1);
% arXiv 1705.09022v4 (secciones 2--4);
% Conway--Norton, pp. 311 y 314--315, tabla 3.
% \bibitem{chenlamshimakura2016}
% H.-Y. Chen, C. H. Lam y H. Shimakura,
% ``$\mathbb Z_3$-orbifold construction of the Moonshine vertex operator
% algebra and some maximal $3$-local subgroups of the Monster'',
% prepublicación, arXiv:1606.05961 (2016), versión 2 (2017).
% \href{https://arxiv.org/abs/1606.05961}{arXiv:1606.05961}.
% \bibitem{abelamyamada2017}
% T. Abe, C. H. Lam y H. Yamada,
% ``A remark on $\mathbb Z_p$-orbifold constructions of the Moonshine
% vertex operator algebra'', prepublicación, arXiv:1705.09022 (2017).
% \href{https://arxiv.org/abs/1705.09022}{arXiv:1705.09022}.
% \bibitem{kmconwaynorton1979}
% J. H. Conway y S. P. Norton, ``Monstrous Moonshine'',
% \emph{Bulletin of the London Mathematical Society} \textbf{11},
% 308--339 (1979).
% \href{https://doi.org/10.1112/blms/11.3.308}{doi:10.1112/blms/11.3.308}.
